% GENERATED FILE. Edit canonical lessons, then run npm run generate:books.
% canonical-source-hash: fnv1a64:f87f5c02d41c8450
% canonical-lessons: PT-C144-calar-se, PT-C144-temer, PT-C144-transportar, PT-C144-matar, PT-C144-unir

\chapter{To fall silent, To fear, To transport, To kill, To unite}
\label{ch:pt-to-fall-silent-to-fear-to-transport-to-kill-to-unite}
\hlchaptermodality{144}

\begin{chapteropening}
\textbf{By the end of this chapter:} \emph{I can say to fall silent, to fear, to transport, to kill and to unite.}

Complete the last lesson of chapter 144: I can say to fall silent, to fear, to transport, to kill and to unite.
\end{chapteropening}

\section[calar-se]{calar-se — to fall silent}

\label{lesson:PT-C144-calar-se}



\begin{quote}
Before the new one: say the Portuguese for to give back, then the Portuguese for to recognise.
\end{quote}

\subsection*{You'll want to know: calar-se}
\textbf{calar-se} — \textquotedblleft{}to fall silent\textquotedblright{}.

\begin{cousinweb}
\textbf{calar} is usually traced to Late Latin \emph{chalare}, \textquotedblleft{}to let down, to slacken\textquotedblright{}, from Greek \emph{chalan} (Lewis \& Short gives the Latin): you let your voice drop. \emph{-se} means \textquotedblleft{}oneself\textquotedblright{}.
\end{cousinweb}

\begin{grammarlens}[title={What you've built}]
One more word for this chapter.
\end{grammarlens}

\subsection*{Your turn}
\begin{itemize}
\raggedright
  \item \emph{Say it:} \emph{calar-se}
  \item \emph{Say it:} \emph{calar-se}, once more
  \item \emph{Say it:} say \emph{calar-se}
  \item \emph{Recall:} say the Portuguese for to give back, then the Portuguese for to recognise, then say \emph{calar-se} again
\end{itemize}

\begin{tcolorbox}[breakable,colback=teal!4,colframe=teal!35!black,title={Before you move on}]
What does calar-se mean? (\textquotedblleft{}To fall silent\textquotedblright{}.) Say it once more.
\end{tcolorbox}
\section[temer]{temer — to fear}

\label{lesson:PT-C144-temer}



\begin{quote}
Before the new one: say the Portuguese for to recognise, then the Portuguese for to fall silent.
\end{quote}

\subsection*{You'll want to know: temer}
\textbf{temer} — \textquotedblleft{}to fear\textquotedblright{}.

\begin{cousinweb}
\textbf{temer} $\leftarrow$ Latin \emph{timere}, \textquotedblleft{}to fear\textquotedblright{} (EtymDB 2.1) — inherited.
\end{cousinweb}

\begin{grammarlens}[title={What you've built}]
One more word for this chapter.
\end{grammarlens}

\subsection*{Your turn}
\begin{itemize}
\raggedright
  \item \emph{Say it:} \emph{temer}
  \item \emph{Say it:} \emph{temer}, once more
  \item \emph{Say it:} say \emph{temer}
  \item \emph{Recall:} say the Portuguese for to recognise, then the Portuguese for to fall silent, then say \emph{temer} again
\end{itemize}

\begin{tcolorbox}[breakable,colback=teal!4,colframe=teal!35!black,title={Before you move on}]
What does temer mean? (\textquotedblleft{}To fear\textquotedblright{}.) Say it once more.
\end{tcolorbox}
\section[transportar]{transportar — to transport}

\label{lesson:PT-C144-transportar}



\begin{quote}
Before the new one: say the Portuguese for to fall silent, then the Portuguese for to fear.
\end{quote}

\subsection*{You'll want to know: transportar}
\textbf{transportar} — \textquotedblleft{}to transport\textquotedblright{}.

\begin{cousinweb}
\textbf{transportar} $\leftarrow$ Latin \emph{transportare}, \textquotedblleft{}to carry across\textquotedblright{} (EtymDB 2.1) — a learned borrowing from Latin.
\end{cousinweb}

\begin{grammarlens}[title={What you've built}]
One more word for this chapter.
\end{grammarlens}

\subsection*{Your turn}
\begin{itemize}
\raggedright
  \item \emph{Say it:} \emph{transportar}
  \item \emph{Say it:} \emph{transportar}, once more
  \item \emph{Say it:} say \emph{transportar}
  \item \emph{Recall:} say the Portuguese for to fall silent, then the Portuguese for to fear, then say \emph{transportar} again
\end{itemize}

\begin{tcolorbox}[breakable,colback=teal!4,colframe=teal!35!black,title={Before you move on}]
What does transportar mean? (\textquotedblleft{}To transport\textquotedblright{}.) Say it once more.
\end{tcolorbox}
\section[matar]{matar — to kill}

\label{lesson:PT-C144-matar}



\begin{quote}
Before the new one: say the Portuguese for to fear, then the Portuguese for to transport.
\end{quote}

\subsection*{You'll want to know: matar}
\textbf{matar} — \textquotedblleft{}to kill\textquotedblright{}.

\begin{cousinweb}
\textbf{matar} is usually traced to Latin \emph{mactare}, \textquotedblleft{}to sacrifice, to slaughter\textquotedblright{} (EtymDB 2.1), though some trace it to other sources; its origin is disputed.
\end{cousinweb}

\begin{grammarlens}[title={What you've built}]
One more word for this chapter.
\end{grammarlens}

\subsection*{Your turn}
\begin{itemize}
\raggedright
  \item \emph{Say it:} \emph{matar}
  \item \emph{Say it:} \emph{matar}, once more
  \item \emph{Say it:} say \emph{matar}
  \item \emph{Recall:} say the Portuguese for to fear, then the Portuguese for to transport, then say \emph{matar} again
\end{itemize}

\begin{tcolorbox}[breakable,colback=teal!4,colframe=teal!35!black,title={Before you move on}]
What does matar mean? (\textquotedblleft{}To kill\textquotedblright{}.) Say it once more.
\end{tcolorbox}
\section[unir]{unir — to unite}

\label{lesson:PT-C144-unir}



\begin{quote}
Before the new one: say the Portuguese for to transport, then the Portuguese for to kill.
\end{quote}

\subsection*{You'll want to know: unir}
\textbf{unir} — \textquotedblleft{}to unite\textquotedblright{}.

\begin{cousinweb}
\textbf{unir} $\leftarrow$ Latin \emph{unire}, \textquotedblleft{}to make one\textquotedblright{} (EtymDB 2.1), from \emph{unus}, \textquotedblleft{}one\textquotedblright{} — a learned borrowing from Latin.
\end{cousinweb}

\begin{grammarlens}[title={What you've built}]
One more word for this chapter.
\end{grammarlens}

\subsection*{Your turn}
\begin{itemize}
\raggedright
  \item \emph{Say it:} \emph{unir}
  \item \emph{Say it:} \emph{unir}, once more
  \item \emph{Say it:} say \emph{unir}
  \item \emph{Recall:} say the Portuguese for to transport, then the Portuguese for to kill, then say \emph{unir} again
\end{itemize}

\begin{tcolorbox}[breakable,colback=teal!4,colframe=teal!35!black,title={Before you move on}]
What does unir mean? (\textquotedblleft{}To unite\textquotedblright{}.) Say it once more.
\end{tcolorbox}
